\ifx\allfiles\undefined
\documentclass[12pt, a4paper, oneside, UTF8]{ctexbook}
\def\path{../config}
\input{\path/package.tex}

% 定理环境设置
\input{\path/theorem.tex}

\input{\path/custom.tex}

\input{\path/cover_package} 

% 封面页设置
\def\myTitle{量子力学：理论与应用
}
\def\mySubtitle{Quantum Mechanics: Theory and Application}
\def\myAuthor{Max}
\def\myCreateDateCover{2026年10月11日}
\def\myUpdateDateCover{ \today}
\def\myDateForeword{}
% 前言页设置
\def\myForeword{前言}
\def\myForewordText{

}
\def\mySubheading{}
\begin{document}
% \input{\path/cover_text.tex}

%%%%%%%% 前言页 %%%%%%%%

\newpage
\thispagestyle{empty}
\begin{center}
    \Huge\textbf{\myForeword}
\end{center}
\myForewordText
\begin{flushright}
    \begin{tabular}{c}
        \myDateForeword
    \end{tabular}
\end{flushright}

% 目录页
\newpage
\pagestyle{plain}
\setcounter{page}{1}
\pagenumbering{Roman}
\tableofcontents

% 正文页
\newpage
\pagenumbering{arabic}
\setcounter{chapter}{-1}
\setcounter{page}{1}

\pagestyle{fancy}
% \fancyhead[L]{量子力学笔记}
% \fancyhead[C]{中间页眉}
% \fancyhead[R]{右侧页眉}
% \fancyfoot[C]{\thepage}
\fancyhf{}
% 设置页眉
\fancyhead[L]{\leftmark}    % 左侧：文章标题（\leftmark 通常是章标题）
\fancyhead[C]{ }                       % 中间：留空（或自定义内容）
\fancyhead[R]{\rightmark}   % 右侧：节标题（\rightmark 通常是节标题）


% 设置页脚（示例：中间显示页码）
\fancyfoot[C]{\thepage}

\renewcommand{\headrulewidth}{0.4pt}
\renewcommand{\footrulewidth}{0pt}








\else
\fi

\chapter{波函数与薛定谔方程}
\begin{examct}{}{}
    波函数及其统计解释、薛定谔方程、态叠加原理、几率守恒定律、定态薛定谔方程
\end{examct}
\begin{examrqm}{}{}
    \begin{enumerate}
        \item 掌握波函数及统计解释、薛定谔方程、态叠加原理、几率守恒定律 
        \item 会灵活运用所学知识解决实际问题
    \end{enumerate}
\end{examrqm}


\section{波函数}

量子力学中求解问题的基本方法是求解\textbf{薛定谔(Schrödinger)方程},得到粒子的波函数 $\Psi(\mathbf{r},t)$ 。
\begin{equation}
    \boxed{\mathrm{i}\hbar\frac{\partial\Psi}{\partial t} = -\frac{\hbar^2}{2m}\nabla^2\Psi + V(\mathbf{r},t)\Psi}
\end{equation}
描述微观粒子波动性的方程, 其中 $\hbar$ 是约化普朗克(Planck)常数, $m$ 是质量, $\nabla^2$ 是 Laplace 算子, $V(\mathbf{r},t)$ 是势函数。
\begin{equation}
    \hbar=\frac{h}{2\pi}=1.054\ 572\times 10^{-34}\,\mathrm{J}\cdot\mathrm{s}
\end{equation}
薛定谔方程的作用和地位从逻辑上将相当于牛顿第二定律,即给定适当的初始条件(一般来说, $\Psi(\mathbf{r},0)$ ),通过薛定谔方程确定以后所有时刻的波函数 $\Psi(\mathbf{r},t)$ 。
\subsection{波函数的统计解释}
关于波函数的物理意义,以波尔为首的哥本哈根学派认为“波函数”表征“概率波”，即波函数的模平方 $\left|\Psi(\mathbf{r},t)\right|^2$ 表示在任意时刻 $t$ 时刻粒子在空间 $\mathbf{r}$ 处的单位体积中出现的概率。
\begin{equation}
    \int_{V}\left|\Psi(\mathbf{r},t)\right|^2\mathrm{d}\mathbf{r} = \{\text{在}t\text{时刻发现粒子处于空间}r\text{处的单位体积中出现的概率}\}
\end{equation}
所以,波函数的统计解释是“概率分布”或“概率密度”。波函数满足
\begin{equation}
    \int_{\infty}\left|\Psi(\mathbf{r},t)\right|^2\mathrm{d}\mathbf{r} = 1
\end{equation}

\subsection{态叠加原理}

\section{概率流守恒定律}

\section{定态薛定谔方程}
考虑 $\hat{H}$ 不依赖于时间的情况, 即
\begin{equation}
    i\hbar\frac{\partial}{\partial t}\Psi(x,t)=\hat{H}\Psi(x,t) = \left[-\frac{\hat{p}^2}{2m}+\hat{V}(x)\right]\Psi(x,t)
\end{equation}
这就是定态薛定谔方程。采用分离变量法。令
\begin{equation}
    \Psi(x,t)=\psi(x)T(t)
\end{equation}
代入得
\begin{equation}
    \boxed{T(t)=e^{-iEt/\hbar}}
\end{equation}
称为\textbf{时间演化因子}。
\begin{equation}
    \boxed{\hat{H}\psi(x,t)=E\psi(x)}\label{eq:TISE}
\end{equation}
该式为 $\hat{x}$ 表象下的能量本征方程。解为能量本征态，即\textbf{“定态”}。\eqref{eq:TISE} 即为\textbf{定态薛定谔方程}。
所以有波函数为
\begin{equation}
    \Psi(x,t)=A\psi(x)e^{-iEt/\hbar}
\end{equation}
其中 $A$ 为归一化因子。

\section{维里定理}
（the virial theorem）

定态下满足,两倍的动能期望值等于位矢量与势能梯度点积的期望值
\begin{equation}
    2\left\langle \hat{T} \right\rangle = \left\langle \vec{r}\cdot \nabla V(\vec{r}) \right\rangle
\end{equation}

\section{F-H定理}
（Feynman-Hellmaon Theorem）
\begin{equation}
    \bra{\psi_n}\frac{\partial }{\partial \lambda}\hat{H}\ket{\psi_n}=\frac{\partial }{\partial \lambda}E_n
\end{equation}
\section{附录}
\ifx\allfiles\undefined
\end{document}
\fi